\documentclass[../../main.tex]{subfiles}
\begin{document}

{\bf [解]}

\begin{figure}[htb]
  \centering
  \begin{tikzpicture}[scale=1.5, >=stealth, every node/.style={font=\small}]

    % --- 1. パラメータ設定 (見やすい比率として a = 2.8 を採用) ---
    % a = 2.8 => y_P = 0.25 * 2.8^2 = 1.96
    % 半径 r = 1 + 1.96 = 2.96
    % 焦点 F(0, 1), 準線接点 H(2.8, -1)
    % 垂線の足 I(0, 1.96), G(0, 1.96 + 0.96) = (0, 2.92)
    \pgfmathsetmacro{\aVal}{2.8}
    \pgfmathsetmacro{\yP}{0.25 * \aVal * \aVal}
    \pgfmathsetmacro{\radR}{1.0 + \yP}
    \pgfmathsetmacro{\yG}{2.0 * \yP - 1.0}

    \coordinate (O) at (0, 0);
    \coordinate (P) at (\aVal, \yP);
    \coordinate (F) at (0, 1.0);
    \coordinate (H) at (\aVal, -1.0);
    \coordinate (I) at (0, \yP);
    \coordinate (G) at (0, \yG);

    % 角度計算 (PH 方向: -90°, PF 方向)
    \pgfmathsetmacro{\angPF}{atan2(1.0 - \yP, -\aVal)} % PF の方向角 (約 -161°)
    \pgfmathsetmacro{\angTheta}{-90.0 - \angPF}         % theta = 角 HPF (約 71°)

    % --- 2. 面積領域の塗りつぶし ---
    % (i) 扇形 PFH (面積 S(a)) : オレンジ透過
    \fill[orange!20, opacity=0.7]
    (P) -- (H) arc[start angle=-90, end angle=\angPF, radius=\radR] -- cycle;

    % (ii) 三角形 PGF (面積 T(a)) : 青色ハッチング
    \fill[pattern=north east lines, pattern color=blue!70, opacity=0.6]
    (P) -- (G) -- (F) -- cycle;

    % --- 3. 座標軸 ---
    \draw[->, thick] (-1.2, 0) -- (4.8, 0) node[right] {$x$};
    \draw[->, thick] (0, -1.8) -- (0, 4.5) node[above] {$y$};
    \node[below left=2pt] at (O) {$O$};

    % --- 4. 準線 y = -1 ---
    \draw[dashed, gray!70, thick] (-1.2, -1.0) -- (4.5, -1.0);
    \node[left=2pt, font=\scriptsize] at (0, -1.0) {$-1$};
    \node[right=2pt, font=\scriptsize, gray!80!black] at (4.5, -1.0) {準線 $y = -1$};

    % --- 5. 放物線 y = 1/4 x^2 ---
    \draw[thick, teal!85!black, domain=-1.0:3.8, samples=60]
    plot (\x, {0.25*\x*\x});
    \node[teal!90!black, font=\footnotesize, right] at (3.4, {0.25*3.4*3.4}) {$y = \frac{1}{4}x^2$};

    % --- 6. 円 C (中心 P, 半径 r) ---
    \draw[very thick, blue!80!black] (P) circle (\radR);
    \node[blue!85!black, font=\normalsize, above right] at ($ (P) + (35:\radR) $) {$C$};

    % --- 7. 線分と直交関係 ---
    % 半径 PH, PF, PG
    \draw[very thick, red!85!black] (P) -- (H) node[midway, right=1pt, font=\scriptsize] {$r$};
    \draw[very thick, blue!85!black] (P) -- (F) node[midway, below=1pt, font=\scriptsize] {$r$};
    \draw[thick, blue!75!black] (P) -- (G);

    % 弦 GF (y 軸上の切断線分)
    \draw[line width=2.0pt, purple!85!black] (G) -- (F);
    \node[purple!90!black, font=\footnotesize, left=2pt] at ($ (G)!0.5!(F) $) {弦 $GF$};

    % 垂線 PI (水平線)
    \draw[thick, dashed, orange!90!black] (P) -- (I);

    % --- 8. 直角記号 ---
    % I における直角記号 (PI ⊥ GF)
    \draw[thin, gray!70]
    ($ (I) + (0.15, 0) $) --
    ($ (I) + (0.15, -0.15) $) --
    ($ (I) + (0, -0.15) $);

    % H における直角記号 (PH ⊥ 準線)
    \draw[thin, gray!70]
    ($ (H) + (-0.15, 0) $) --
    ($ (H) + (-0.15, 0.15) $) --
    ($ (H) + (0, 0.15) $);

    % --- 9. 角度の弧とラベル ---
    % 角 HPF = theta
    \draw[thick, purple!85!black]
    ($ (P) + (0, -0.7) $) arc[start angle=-90, end angle=\angPF, radius=0.7];
    \node[purple!90!black, font=\footnotesize] at ($ (P) + (-105:0.95) $) {$\theta$};

    % 角 FPI = pi/2 - theta
    \draw[thick, gray!70]
    ($ (P) + (\angPF:0.55) $) arc[start angle=\angPF, end angle=-180, radius=0.55];
    \node[gray!80!black, font=\scriptsize] at ($ (P) + (-172:0.9) $) {$\frac{\pi}{2}-\theta$};

    % --- 10. 主要点のプロットとラベル ---
    \fill (O) circle (0.6pt);
    \fill[teal!85!black] (F) circle (0.8pt) node[left=2pt, teal!90!black] {$F(0,1)$};
    \fill[red!85!black]  (H) circle (0.8pt) node[below=2pt, red!90!black] {$H$};
    \fill[blue!85!black] (P) circle (0.9pt) node[right=1pt, blue!90!black] {$P\left(a, \frac{1}{4}a^2\right)$};
    \fill[orange!90!black] (I) circle (0.7pt) node[below left=1pt, orange!90!black] {$I$};
    \fill[purple!85!black] (G) circle (0.8pt) node[above left=1pt, purple!90!black] {$G$};

  \end{tikzpicture}
  \caption{放物線 $y=\dfrac{1}{4}x^2$，準線に接する円 $C$，弦 $GF$ および扇形 $PFH$}
  \label{fig:parabola_circle_chord}
\end{figure}

$y=-1$を準線，$F(0,1)$を焦点とする放物線は
\begin{align}
  y = \dfrac{1}{4} x^2
\end{align}
である．
この放物線上の点 $(a, \dfrac{1}{4} a^2)$ を中心とし，$y = -1$ に接する，つまり半径 $r$ として $r = 1 + \dfrac{1}{4} a^2$ の円 $C$ は
\begin{align}
  C : (x - a)^2 + \Bigl(y - \dfrac{1}{4} a^2\Bigr)^2 = \Bigl(1 + \dfrac{1}{4} a^2\Bigr)^2
\end{align}
である．

$GF$ の中点 $I$ とすると，$PI \perp GF$ になっていることに注意して，$\triangle PGF$の面積$T(a)$は
\begin{align}
  T(a) = \dfrac{1}{2} \cdot |IP| \cdot |GF| = \dfrac{1}{2} \cdot |IP| \cdot (2 |FI|) = |IP||FI| \label{eq:1}
\end{align}
である．
$PH$ と $PF$ のなす角 $\theta$ とすると，$\angle IPF = \dfrac{\pi}{2} - \theta$だから$|IP|=r\cos\left(\dfrac{\pi}{2} - \theta\right), |FI|=r\sin\left(\dfrac{\pi}{2} - \theta\right)$だから\cref{eq:1}に代入して
\begin{align}
  T(a)
  &= r \sin\Bigl(\dfrac{\pi}{2} - \theta\Bigr) r \cos\Bigl(\dfrac{\pi}{2} - \theta\Bigr) \\
  &= r^2 \sin\theta \cos\theta \label{eq:2}
\end{align}
である．
一方で扇形$PFH$の面積$S(a)$は
\begin{align}
  S(a) = \dfrac{1}{2} r^2 \theta \label{eq:3}
\end{align}
である．よってこれらの比は
\begin{align}
  \dfrac{T(a)}{S(a)}
  &= \dfrac{r^2 \sin\theta \cos\theta}{\dfrac{1}{2} r^2 \theta} \\
  &= \dfrac{2 \sin\theta \cos\theta}{\theta} \\
  &= \dfrac{\sin 2\theta}{\theta}\label{eq:4}
\end{align}
である．以下$a\to\infty$でのこの極限を求めよう．

さて，ここで
$\overrightarrow{PH} =
\begin{pmatrix} 0 \\ -r
\end{pmatrix}$， $\overrightarrow{PF} =
\begin{pmatrix} -a \\ 1 - \dfrac{1}{4} a^2
\end{pmatrix}$
より，ベクトルの内積の公式より
\begin{align}
  \cos\theta
  &= \dfrac{\overrightarrow{PH} \cdot \overrightarrow{PF}}{|\overrightarrow{PH}| |\overrightarrow{PF}|} = \dfrac{-r \Bigl(1 - \dfrac{1}{4} a^2\Bigr)}{r^2} \\
  &= \dfrac{-1 + \dfrac{1}{4} a^2}{r} = \dfrac{-4 + a^2}{4 + a^2} \label{eq:5}
\end{align}
である．
よって$a \to \infty$ の時，$\cos\theta = \dfrac{-4/a^2 + 1}{4/a^2 + 1} \to 1$ より $\theta \to 0$ であるので，\cref{eq:4}に代入して求める極限値は
\begin{align}
  \lim_{a\to\infty}\dfrac{T(a)}{S(a)}
  = \lim_{\theta\to 0}\dfrac{\sin 2\theta}{2\theta} \cdot 2
  =  2
\end{align}
である．

\end{document}
